\section{Exact partial cancellation with cycle-block kernels}
\label{sec:cycle-block-cancellation}

The all-forest obstructions of Section~\ref{sec:representative-activation}
leave a different local question: how much of a stalled core can be
cancelled, with what certified residual cost?  Report~116 supplies a
weighted matrix optimizer parameterized by the largest individual block
cycle rank.  We adapt that optimizer in
\path{fastunknot/cycle_block_cancellation.py}, add a polynomial exact
single-deferral test, replace recursive enumeration by an explicit stack,
and reject cyclic prefixes using rollback reachability.  This is an
additive matrix API.  It is not inserted into the recognizer's verdict
path, because a partial plan still requires ordered geometric transport
and residual processing.

\subsection{Positive cancellation and source-private plans}

Let $A$ be the sparse nonnegative matrix of original unsigned occurrences,
with a nonnegative binary weight $w_r$ per row.  A positive unit pivot
$(r,c)$ removes its row and column and sends every surviving entry to
$A'_{sh}=A_{sh}+A_{sc}A_{rh}$.  A source-private order
$(r_i,c_i)_{i=1}^k$ uses distinct rows and columns with
$A_{r_i c_i}=1$ and $A_{r_j c_i}=0$ for $j>i$.  Its selected donors
therefore remain unchanged by earlier pivots; passive rows may acquire
fill.  The objective is to maximize the total weight of selected rows.

For unrestricted positive pivots, a set of rows is eliminable exactly
when its original submatrix has a source-private order.  Indeed, failure
of private peeling leaves a nonempty stopping set $Q$ whose column
totals are all zero or at least two.  If a pivot row is outside $Q$, its
new totals are $s'_h=s_h+s_c A_{rh}$.  If it belongs to $Q$, the surviving
set is nonempty and has totals $s'_h=s_h+(s_c-2)A_{rh}$.  Neither formula
can create a unit total.  Thus positive fill cannot remove the entire
stopping set.  Conversely a private order is a valid positive sequence.
This is the inherited positive/private equivalence, not a new geometric
cancellation theorem.  Restrictions to designated original unit pairs
are supported for private orders only.

A globally private eligible pair can be forced without reducing the
optimum: prepend it to any plan omitting its row, or move its row to the
front of a plan containing it.  Its column vanishes on every other row.
Nonnegative weights are needed here.  The implementation first runs this
queue to exhaustion, retaining its exact prefix.

\begin{proposition}[Exact minimum-weight single deferral]
Suppose forced private peeling has stalled on a nonempty residual row
set $R$.  Put $\mu=\min_{r\in R}w_r$ and $W=\sum_{r\in R}w_r$.
If deleting a row $r$ of weight $\mu$ permits complete private peeling
of $R\setminus\{r\}$, the residual optimum is $W-\mu$.
\end{proposition}
\begin{proof}
Complete private elimination of $R$ is impossible: a valid complete
order would begin with a globally private eligible column, contradicting
the stalled queue.  Every feasible selected subset therefore omits at
least one row and loses weight at least $\mu$.  The proposed trace attains
the upper bound $W-\mu$.  This also covers $\mu=0$.  The already forced
prefix is compatible with that trace and adds its fixed weight.
\end{proof}

We try minimum-weight rows in label order and stop on the first successful
trace.  The default allowance is eight trials; its exhaustion, like failure,
falls through to the complete block optimizer.  The allowance can be set to
\code{None}.  If $m$
is the number of residual rows, $z$ the nonzero count and $n$ the number
of active columns, the prepass uses at most $m$ sparse peeling queries,
each costing $O(m+z+n\log(n+2))$ operations in the straightforward
implementation.  Thus its overhead is polynomial in the encoded input.
The option \code{single\_deferral=False} disables it for controlled
block-kernel experiments.  It runs only after exhausted forced peeling;
an already completely peelable input never loses a row.

\subsection{From matchings to six-state blocks}

Form the bipartite support graph on rows and columns having an eligible
unit entry.  Retain every nonzero edge in those columns, including heavy
entries; only eligible unit edges may be matched.  Orient matched edges
row-to-column and all other edges column-to-row.  Directed cycles are
exactly alternating cycles.  Their absence is equivalent to a private
order of the matched pairs by reverse topological sorting.  This is the
classical uniquely restricted matching formulation; related
parameterized algorithms include \cite{cycle-urm}.

Every simple cycle lies in a single vertex-biconnected block.  Blocks
therefore communicate only through the matching capacity of each shared
articulation vertex.  A block sends two scores according to whether it
occupies its parent vertex.  At a vertex, either no child occupies it or
one child does.  Summing the unoccupied child responses and taking the
largest nonnegative occupied-minus-unoccupied gain gives the two exact
responses to the parent block.  Witnesses are stored in a shared DAG.

Inside one block, suppress vertices of degree two, keeping the parent
articulation.  A path exposes its first and last matching flags
$(\alpha,\beta)$ and a direction $d\in\{-,0,+\}$: nonzero means that
every oriented edge points consistently along the whole path.  There
are at most four zero-direction states and one state in each nonzero
direction, since a consistent direction uniquely determines every
matching flag.  A constant-state dynamic program along the path includes
the appropriate occupied/free descendant score at each internal vertex.

For block cycle rank $b=|E|-|V|+1\ge2$, the identity
$\sum_v(\deg(v)-2)=2b-2$ gives at most $2b-2$ branch vertices.  Keeping
the parent adds at most one more; suppression then gives at most $3b-2$
paths.  Rank-one blocks are cycles opened as one root-to-root path.
Enumerate the six-state choices, reject conflicting endpoint capacities,
and require the directed path skeleton to be acyclic.  A directed cycle
in the original block traverses entire degree-two paths, so this skeleton
criterion is exact, including root-to-root self-loops.

Writing $\kappa$ for the maximum individual block rank, this yields
$\operatorname{poly}(\Lambda)2^{O(\kappa)}$ bit complexity.  Scores sum
binary input row weights and have polynomial bit length.  The number of
blocks contributes a polynomial factor, not an exponent depending on
the sum of their ranks.  Cactus tractability is classical; the delivered
refinement combines weighted eligibility, six-state paths and the
individual-block parameter.

\subsection{Cycle rejection before the leaf and unbounded stack depth}

The reference enumerator checks directed acyclicity only at complete
assignments.  Our maintained implementation stores each terminal's
reachability set as a bit mask.  Adding $u\to v$ creates a cycle exactly
when $u=v$ or $v$ already reaches $u$.  Otherwise, for every ancestor of
$u$, including $u$, insert $v$ and every descendant of $v$ into its mask.
Save only changed masks, then restore them on backtracking.

This update is the exact new transitive closure: every new path consists
of an old path to $u$, the new arc, and an old path from $v$.  Since the
prefix was acyclic, no later choice can repair a rejected cycle.  Pruning
therefore loses no feasible matching.  Closure maintenance contributes
only polynomial factors in the block size and preserves the
$2^{O(\kappa)}$ bound.  Its overhead can still hurt small instances.

An explicit stack holds choice iterators, accumulated scores and rollback
records.  This removes the delivered implementation's hard maximum of
256 paths.  The optional operational path cap defaults to 128 but accepts
\code{None}; setting every cap to \code{None} now exposes the finite
uncapped algorithm without a recursion-depth qualification.  A completed
270-path test on $K_{2,270}$ disables the single-deferral prepass.  Its
optimum is one: two matched row/column pairs would alternate on a
four-cycle.

\begin{proposition}[A family with a lower search exponent]
On the all-unit two-row, $s$-column matrix with unit row weights and
$s\ge3$, the delivered
leaf-only skeleton enumeration takes $\Theta(s^3)$ search operations,
whereas the maintained cycle-prefix enumeration takes $\Theta(s^2)$.
With single deferral enabled, peeling settles the matrix in $O(s)$
arithmetic work, apart from external-label sorting.
\end{proposition}
\begin{proof}
The support graph is $K_{2,s}$.  Its skeleton has two row terminals and
$s$ length-two paths.  Each path permits exactly three states: neither
endpoint matched, only the first matched, or only the second matched.
At a prefix of $i$ paths, endpoint capacity alone permits
$1+2i+i(i-1)=i^2+i+1$ assignments.  The last term consists of two matched
rows on different paths; those assignments form directed two-cycles in
the skeleton.  The reference search retains them until all $s$ paths are
assigned.  Summing its prefix counts gives $\Theta(s^3)$ work, even
though only $s^2+s+1$ leaves reach its final acyclicity check.

Immediate cycle rejection retains just $1+2i$ assignments at depth $i$.
Each has at most three extensions and closure updates on two terminals
cost constant work.  The sum is $\Theta(s^2)$, with $2s+1$ complete
assignments.  Finally, forced peeling initially stalls; omitting either
row lets the remaining row pivot on any column.  The minimum-weight
single-deferral test therefore proves the optimum one after two linear
incidence scans.  Sorting arbitrary external labels is a separate cost.
\end{proof}

Callbacks run even on empty input.  Caller exceptions propagate,
including exceptions of the public \code{ResourceLimit} type;
they are distinguished from internal budget exhaustion.  This corrects
the delivered implementation's broad catch, which could turn that caller
exception into a resource-status result.  A resource refusal remains
inconclusive.  The work counter is not a hard memory or wall-clock bound.

\subsection{Deferred forests and the geometric limit}

If each uncapped row is a mod-two cycle in the zero/one-handle graph, the
pivot columns of a partial private plan can be removed without disconnecting
any original component.  At step $i$, its donor row avoids all earlier
pivots and uses $c_i$ exactly once.  The surviving parity cycle witnesses
that $c_i$ is not a bridge.  A spanning forest disjoint from all selected
pivots therefore exists, and any prescribed disjoint forest extends to one.
Restoring or removing other forest columns does not affect the private
trace.  Consequently optimizing partial cancellation on all columns
already optimizes over all lower spanning forests.  This extends the
earlier full-cancellation argument to weighted partial plans; saturated
counts cannot replace uncapped parity here.

The delivered numerical transport circuit writes donor effects backwards
as $T_i=R_i+\sum_{j>i}A_{r_i c_j}T_j$ and passive effects as
$P_p=R_p+\sum_i A_{p c_i}T_i$.  There are linearly many source coefficient
occurrences.  Scalar linear observables evaluate without materializing
quadratic passive fill, with polynomial bit lengths by the finite
acyclic-path expansion.  This circuit still lacks cyclic order, signs,
ribbon data and protected boundary patterns.  It is retained as delivered
research, not promoted as a geometric successor representation.

\subsection{Native evidence and comparison scope}

The native audit authenticates 126 source parallelity packets, independently
replays the existing maximum-weight top reduction, and reconstructs exact
full-column counts.  All 126 maintained matrix optimizations complete and
agree in objective with report~110's ordered-interface optimizer and its
independent feasible-plan replay.  Of these, 124 completely peel.  The
other two attain the known optimal residual weight one:
\begin{center}
\begin{tabular}{lrrrr}
source & rows & forced & residual & new trials\\
\hline
\path{nested_1_4_6/interior_links}&42&32&10&1\\
\path{interior_loop/interior_links}&9&5&4&1
\end{tabular}
\end{center}
Both are settled by one minimum-weight deferral trial.  The selected
original row is zero in both new plans; the nested case's ordered-interface
plan instead defers row 24.  These are different valid matrix plans, so
downstream basis-dependent data cannot be copied between them.

The direct block implementation before this prepass exhausts two million
events on the nested residual, whose block rank is 11 and whose skeleton
has 20 paths.  That failed audit is retained.  The delivered block solver
likewise reaches its default work cap on that input.  Its current ordered
decomposition has width four and peak 3840 states.  The final prepass uses
544 charged events on the nested matrix and 114 on the loop matrix;
these event counts have different definitions from ordered DP states.
They are not timing ratios.  Native source authentication and compatible
forest construction are executed; geometric partial cancellation is not.

The 50 delivered tests and their 31,581 exhaustive matrix-case roles
(63,162 plans with prepeeling on/off) pass both on the original scratch
copy and with the maintained optimizer substituted.  The randomized
audit additionally checks weights, restrictions, relabelling, positive
transport and 4800 forest choices.  Ten maintained tests add an independent
512-input binary enumeration, 500 weighted restricted inputs, the deep
stack example, malformed packets, empty rows, exact single deferral and
callback/resource separation and prepass fallback.  The supplied large algebra audit also
passes on the original copy.  These checks do not validate geometric
partial transport or prove a global unknot complexity bound.

\begin{center}
\begin{tabular}{lrrrrr}
query & ordered ms & delivered ms & new ms & ordered/new & A/A\\
\hline
Pachner mixed & 0.267 & 0.289 & 0.292 & 0.92 & 0.998\\
nested links & 39.072 & cap & 0.864 & 45.32 & 1.060\\
loop links & 1.987 & 4.273 & 0.259 & 7.67 & 1.031\\
cactus 128 & 337.336 & 15.265 & 15.183 & 21.87 & 1.009\\
$K_{2,48}$ & 7.040 & 66.346 & 0.335 & 21.43 & 1.014\\
\end{tabular}
\end{center}


Here ``ordered/new'' is the median of paired query ratios, not the ratio
of separately computed medians.  The native nested and loop matrix queries
give paired factors about 45.3 and 7.7 against the specified ordered
implementation.  The already peelable Pachner case is slightly slower
than that comparator.  The first uncapped single-deferral prepass run
regressed on the 128-block cactus (48.31\,ms versus the delivered solver's
15.27\,ms).  After limiting the prepass to eight candidates, the final
run is near parity with the delivered solver (15.18 versus 15.26\,ms).
Both runs and their exact tested source snapshots are retained; samples
are not pooled or replaced.

With single deferral disabled, a separate complete-query ablation on
$K_{2,s}$ for $s=16,48,96$ gives median paired delivered/new ratios
2.25, 7.32 and 16.42.  This isolates the iterative cycle-prefix implementation
from the new optimality prepass.  The largest query falls from about
448.27\,ms to 27.08\,ms.  These graphs are synthetic, not claimed native
normal-cut families.

The timings are complete \emph{matrix} queries from already prepared
counts: normalization, discovery, feasible-plan replay and JSON
serialization.  Source geometry, attaching-curve transport and residual
recognition are excluded.  Nine randomized-order rounds and three A/A
controls are retained, together with a separate prefix-pruning ablation.
A capped refusal has no speedup ratio against a completed answer.  The
results compare specified implementations, not all algorithms for uniquely
restricted matching.  No new knot verdict or whole-cut positive follows
from this matrix optimization.  General quasi-polynomial recognition still
requires useful-surface coverage, controlled marked geometric successors
and aggregate hierarchy accounting.
